% A.79 reserved by the root.
\subsection{A valuation minimum in the opposite regular coordinate}

\begin{statement}{A.79}{Proved}
{A logarithmic valuation criterion with explicit additional infinite families}
Let $\lambda>1$ be an even integer, put $k=v_2(\lambda)$,
and factor its odd neighbour as
\[
               \lambda-1=\prod_{q\in Q}q^{e_q}.
\tag{A.79.1}
\]
Define the set of potentially multiplicative contact-zero primes
\[
 P=\{p\text{ odd}:p\mid\lambda\}
                       \cup\bigl(\{2\}\text{ if }k\ge5\bigr).
\tag{A.79.2}
\]
For $p\in P$ and $q\in Q$ define the explicitly computable
positive integer
\[
 w_{p,q}=
 \begin{cases}
 v_p(e_q)+v_p(q^{p-1}-1),&p\text{ odd},\\
 v_2(e_q)+v_2(q^2-1)-1,&p=2.
 \end{cases}
\tag{A.79.3}
\]
Suppose that for every $p\in P$ the minimum of the numbers
$w_{p,q}$, $q\in Q$, is attained at exactly one $q$.
Every noncuspidal Selmer section in A.14 with labeled integral
realization $E_\lambda$ is the section of $\lambda$ in either
of the following proved ranges:
\begin{itemize}
\item $k\notin\{2,4\}$, with an arbitrary labeled integral
comparison and no multiplier restriction;
\item $k\in\{2,4\}$, with its full congruence image conjugate
to that of $s_\lambda$ by $H$.
\end{itemize}
In particular, this recognizes the two infinite families
\[
 \begin{array}{c|c}
 \lambda=191^{2j+1}+1,\quad j\ge0
       &\text{arbitrary labeled integral comparison},\\
 \lambda=3^{2j+1}+1,\quad j\ge2
       &\text{full }H\text{ conjugacy}.
 \end{array}
\tag{A.79.4}
\]
It also recognizes $\lambda=88$ and $120$ with an arbitrary
labeled integral comparison, although their odd neighbours
have two distinct prime factors.
\end{statement}
\begin{proof}
At odd primes, the labeled integral cyclotomic-line calculation
in A.41 assigns contact zero to divisors of $\lambda$, contact
one to divisors of $\lambda-1$, and no contact at other primes.
At two, potentially good reduction excludes a local cusp by
A.37.  In the potentially multiplicative range $k\ge5$,
the exact closest-pair calculation (A.78.8)--(A.78.9) assigns
label zero to any local cusp.  This label calculation uses only
the labeled integral comparison, not its determinant.
Thus the only possible cuspidal labels are zero and one,
and every label-zero cusp has its prime in $P$.

We first show that the global $B$ coordinate for $1-t$ has
no valuation at two.  When $k\ge5$, an ordinary local section
has $x_2=\lambda$ by (A.78.1), and a cuspidal local section
has label zero as just proved.  Both give $v_2(B)=0$.
When $k=1$ or $3$, the mod-two Kummer classes in A.16 give
odd valuation parity for $A_2$ and even parity for $B_2$.
For an ordinary point the alternatives are
$(v_2(x_2),v_2(1-x_2))=(h,0),(0,h),(-h,-h)$ with $h>0$.
The indicated parities force the first alternative, so again
$v_2(B)=0$.  There is no local cusp because $k\le4$ makes
the realized curve potentially good.
Finally, for $k=2$ or $4$ the stated full-$H$ hypothesis gives
the two separate mod-eight Kummer characters.
The twelve potentially good pairs (A.78.10) are distinct
modulo eight; their pair is therefore exactly $(k,0)$.
This proves the required dyadic assertion in precisely the
two ranges of the statement.

At every $q\in Q$, both an ordinary local point and a local
cusp have label one.  At ordinary primes the exact-value
statement (A.78.1) gives the honest valuation $e_q$ for $B$.
At a cusp prime $q$, choose a common finite splitting
extension and use the full integral Tate comparison at every
coefficient $\ell$.  Its triangular diagonal entries are
integral units.  Since $q$ is odd, the cusp Tate parameter
has class $B^2/256$ and valuation $2v_q(B)$; the realized
Tate parameter has valuation $2e_q$.
Canceling the common ramification index gives
\[
        \beta_{q,\ell}=e_q\,z_{q,\ell},
                  \qquad z_{q,\ell}\in\Z_\ell^\times,
        \qquad\beta_q=v_q(B)\in\widehat\Z.
\tag{A.79.5}
\]
At an ordinary $q$ the same equation holds with $z_{q,\ell}=1$.
No square condition is used: $z_{q,\ell}$ is the actual
ratio of the two diagonal entries of the integral comparison.
The mod-two global class fixes the negative sign of $B$,
and all its other prime valuations vanish.  Hence the exact
global Kummer expansion is
\[
                  B=-\prod_{q\in Q}q^{\beta_q},
          \qquad v_\ell(\beta_{q,\ell})=v_\ell(e_q)
                    \quad\text{for every }q,\ell.
\tag{A.79.6}
\]
The product is finite and uses the sign/valuation injection
into $\widehat{\Q^\times}$.

Suppose a local cusp has label zero and prime $p\in P$.
Its regular coordinate gives $B_p=1$ in
$\widehat{\Q_p^\times}$.
Every $q\in Q$ differs from $p$.
The continuous logarithm on local units kills the sign and
finite roots of unity and sends this equality to
\[
                       \sum_{q\in Q}
                                 \beta_{q,p}\log_p(q)=0.
\tag{A.79.7}
\]
This is the actual coefficient-$p$ component of localization:
the principal units form a pro-$p$ group, and the power map
from $\widehat\Z$ to them factors through $\Z_p$.

For odd $p$, Fermat's theorem puts $q^{p-1}$ in
$1+p\Z_p$.  The logarithm is an isometry from this group
to $p\Z_p$, and $p-1$ is a unit, so
\[
                    v_p(\log_p q)=v_p(q^{p-1}-1).
\tag{A.79.8}
\]
For $p=2$, one has $q^2\in1+8\Z_2$ and
$2\log_2(q)=\log_2(q^2)$, whence
\[
                    v_2(\log_2 q)=v_2(q^2-1)-1.
\tag{A.79.9}
\]
These logarithms are nonzero: a positive rational prime
cannot be a local root of unity, since the corresponding
finite power equation would hold in $\Q$ itself.
Together with (A.79.5), formulas (A.79.8)--(A.79.9) prove
that the valuations of the summands in (A.79.7) are
exactly the fixed integers $w_{p,q}$.
In a nonarchimedean field, a finite sum with a unique
summand of smallest valuation has that same finite valuation;
divide by that summand and all other terms lie in $p\Z_p$.
It cannot be zero.  The hypothesis therefore excludes a
label-zero cusp at every $p\in P$.

Every remaining local cusp has label one.  The one-label
recognition theorem A.37 makes $s$ a rational point section.
The primitive labeled isogeny argument of A.54 identifies
its parameter with the given $\lambda$ exactly.
This proves the criterion with the scopes stated.

For (A.79.4), both $191$ and $3$ are primes; primality of
$191$ follows by checking $2,3,5,7,11,13$.
An odd power satisfies
$v_2(q^{2j+1}+1)=v_2(q+1)$, since the alternating quotient
by $q+1$ has an odd number of odd summands and is odd.
Thus the two families have $k=6$ and $k=2$, respectively.
Their sets $Q$ each contain just one prime, so every
minimum in (A.79.3) is unique.  This proves both infinite
families with their stated determinant scopes.

The two composite-neighbour examples exhibit the additional
valuation mechanism explicitly.  For $88=8\cdot11$,
$Q=\{3,29\}$ and $P=\{11\}$.  One has
\[
             v_{11}(3^{10}-1)=2,\qquad
             v_{11}(29^{10}-1)=1.
\tag{A.79.10}
\]
For $120=8\cdot3\cdot5$, $Q=\{7,17\}$ and
$P=\{3,5\}$, with exact weight table
\[
 \begin{array}{c|cc}
             &q=7&q=17\\ \hline
 p=3&1&2\\
 p=5&2&1.
 \end{array}
\tag{A.79.11}
\]
Indeed $7^2-1=48$, $17^2-1=288$, $7^4-1=2400$
and $17^4-1=83520$ have the indicated valuations.
Here $k=3$, so no $H$ hypothesis is needed in either example.
\end{proof}

\paragraph{Range comparison and finite certificate.}
A.41 requires a prime-power odd neighbour and an odd
dyadic valuation.  It does not apply to either family
in (A.79.4), or to the composite-neighbour examples.
For the second family $j\ge2$ gives $\lambda\ge244$,
and the first begins at $192$, so they lie beyond the
finite list of A.45.  Neither coordinate in these families
is a signed power of two: $\lambda-1$ is an odd prime power,
while $\lambda$ has dyadic valuation exactly six or two
and exceeds $64$ or $4$, respectively.
The six actual cusp-permutation coordinate functions are
$\lambda$, $1-\lambda$, $1/\lambda$, $1/(1-\lambda)$,
$\lambda/(\lambda-1)$ and $(\lambda-1)/\lambda$.
For an even integer $\lambda>2$, the first and third are
signed dyadic powers exactly when $\lambda$ is a power
of two; the second and fourth have this property exactly
when $\lambda-1$ is a power of two.
The fifth, a reduced fraction greater than one, can be a
dyadic power only if its denominator is one; the sixth,
a reduced fraction less than one, only if its numerator
is one.  Both alternatives force $\lambda=2$.
Thus A.63 does not already supply these families.
The same direct checks distinguish $88$ and $120$ from
A.41, A.45 and A.63.

The program \path{code/legendre_log_minimum_certificate.py}
records the exact criterion on even parameters through
$5000$, its comparison with these earlier criteria, and
the displayed examples and initial members of both
infinite families.  Its output is retained in
\path{code/legendre_log_minimum_certificate_output.txt}.
The general criterion and the infinite families are
proved above, independently of that finite search.
